\documentclass[11pt,reqno]{amsart}

%==============================================================
% A Diagonal Carre-du-Champ Seed for Sticky Brownian Flows
% --------------------------------------------------------------
% Tier alpha input-existence paper for the construction of
%   Fontes, From Carre-du-Champ Hierarchies to Primitive
%   Divergence  [GammaHierarchyDivergence]
% Companion to:
%   Fontes, Primitive Divergence and Carre-du-Champ Seeds over
%   Nonclassical Noises  [PrimitiveDivergence]
% and to the IBM trilogy
%   [EnergyFirst], [FCDC], [BidegreeAtlas].
%
% Scope: extract the diagonal-supported carre-du-champ seed
% Gamma^(2)_theta from the Le Jan--Raimond two-point sticky
% generator under a stated convention, and verify the
% (Sym),(Vanish),(MatPos),(WMatPos),(FormClos) axioms required
% by GammaHierarchyDivergence. Tier beta (lifting to cylindrical
% functionals of the full flow noise) is left for subsequent work.
%==============================================================

\usepackage[utf8]{inputenc}
\usepackage[T1]{fontenc}
\usepackage{lmodern}
\usepackage[margin=1.15in]{geometry}
\usepackage{amsmath,amssymb,amsthm,mathtools}
\usepackage{mathrsfs}
\usepackage{microtype}
\usepackage{hyperref}
\usepackage{cleveref}

\emergencystretch=3em

% --- theorem environments (compatible with companion papers)
\theoremstyle{plain}
\newtheorem{theorem}{Theorem}[section]
\newtheorem{proposition}[theorem]{Proposition}
\newtheorem{lemma}[theorem]{Lemma}
\newtheorem{corollary}[theorem]{Corollary}

\theoremstyle{definition}
\newtheorem{definition}[theorem]{Definition}
\newtheorem{example}[theorem]{Example}
\newtheorem{assumption}[theorem]{Assumption}
\newtheorem{convention}[theorem]{Convention}

\theoremstyle{remark}
\newtheorem{remark}[theorem]{Remark}

% --- macros
\newcommand{\R}{\mathbb{R}}
\newcommand{\N}{\mathbb{N}}
\newcommand{\E}{\mathbb{E}}
\newcommand{\Prob}{\mathbb{P}}
\newcommand{\Dom}{\mathrm{Dom}}
\newcommand{\cA}{\mathcal A}
\newcommand{\cE}{\mathcal E}
\newcommand{\cN}{\mathcal N}
\newcommand{\cT}{\mathcal T}
\newcommand{\cU}{\mathcal U}
\newcommand{\Dth}{D_{\theta}}
\newcommand{\divth}{\delta_{\theta}}

\title{A Diagonal Carr\'e-du-Champ Seed\\
       for Sticky Brownian Flows}
\author{Ramiro Fontes}
\date{\today}

\begin{document}

\begin{abstract}
The companion construction~\cite{GammaHierarchyDivergence}
showed that, given a matrix-positive, weighted-matrix-positive,
form-closable carr\'e-du-champ seed
$\Gamma^{(2)}\colon\cA\times\cA\to L^{1}(\Omega)$ on a unital
test algebra $\cA\subset L^{\infty}(\Omega)\cap L^{2}(\Omega)$
dense in $L^{2}$, one constructs a Hilbert module of one-forms
$\cU_{\Gamma}$, a closed operator derivative
$D_{\cN}=\overline{D_{0}}$, and a closed primitive divergence
$\delta_{\cN}=D_{\cN}^{*}$ satisfying the integration-by-parts
identity. That construction is universal in its inputs but
conditional on existence of the seed.

The present paper supplies a first concrete input. Working under
the Le Jan--Raimond two-point generator/domain convention for the
Brownian sticky flow, with sticky parameter $\theta\in[0,\infty]$,
diagonal interaction coefficient $c_{\theta}\ge 0$, and diagonal
reference measure $\mu_{\Delta}$ (all three tied to the cited
convention; see \cref{conv:LJR} below), we extract from the
two-point sticky generator $L^{(2)}_{\theta}$ a diagonal
carr\'e-du-champ seed
\[
  \Gamma^{(2)}_{\theta}(f,g)
  \;=\;\frac{c_{\theta}}{2}\,f'(x)g'(x)\,\mu_{\Delta}(dx)
\]
on a smooth test class $\cT^{(2)}$ on the diagonal
$\Delta\subset\R^{2}$. We verify the five axioms required
by~\cite{GammaHierarchyDivergence}: symmetry \textsc{(Sym)},
constancy vanishing \textsc{(Vanish)}, matrix positivity
\textsc{(MatPos)}, weighted matrix positivity \textsc{(WMatPos)},
and form closability \textsc{(FormClos)}. The first four are
immediate from the square-of-sum identity
$\sum_{i,j}U_{i}U_{j}f_{i}'(x)f_{j}'(x)=\bigl(\sum_{i}U_{i}f_{i}'(x)\bigr)^{2}\ge 0$;
the fifth is the standard form-closability of the
$\mu_{\Delta}$-weighted Dirichlet form $\frac{c_{\theta}}{2}\int(f')^{2}\,d\mu_{\Delta}$
under the stated convention.

By the transfer theorem of~\cite{GammaHierarchyDivergence}, this
seed yields a Hilbert one-form module $\cU_{\Gamma_{\theta}}$, a
closed operator derivative $\Dth=\overline{D_{0}^{\theta}}$ on
$L^{2}(\R,\mu_{\Delta})$, and a closed densely-defined primitive
divergence $\divth=\Dth^{*}$ with the integration-by-parts
identity
$\E[F\,\divth(u)]=\langle\Dth F,u\rangle_{\cU_{\Gamma_{\theta}}}$.
The pair $(\Dth,\divth)$ lives on the diagonal-supported space
$L^{2}(\R,\mu_{\Delta})$, not yet on cylindrical functionals of
the full flow noise of $\Phi_{\theta}$.

This is a Tier $\alpha$ result. The lift to cylindrical
functionals of the LJR flow noise --- which would produce the
operator-derivative pair on $L^{2}(\cN_{\theta})$ rather than on
$L^{2}(\R,\mu_{\Delta})$ --- is Tier $\beta$ and is left for
subsequent work~(\cref{sec:tier-beta}). We do not claim a
construction over Brownian web, percolation noise, or any other
black noise.
\end{abstract}

\maketitle

%==============================================================
\section{Introduction}\label{sec:intro}
%==============================================================

\subsection{Where this paper sits}\label{sec:intro-where}

The paper~\cite{PrimitiveDivergence} formulated the
operator-derivative problem over a nonclassical noise and
identified its primitive datum as a tuple
$(\{\Gamma^{(N)}_{\cN}\}_{N\ge 2},\delta_{\cN},D_{\cN})$, with
the bilinear two-point carr\'e-du-champ
$\Gamma_{\cN}^{(2)}=\Gamma^{(2)}_{\cN,\{i,j\}}$ as the primary
candidate seed and
$\{\Gamma^{(N)}_{\cN}\}_{N\ge 3}$ as compatibility data. The
companion paper~\cite{GammaHierarchyDivergence} closed the
construction direction, end-to-end:
\[
  \Gamma^{(2)}_{\cN}
  \;\longrightarrow\;
  \cU_{\Gamma}
  \;\longrightarrow\;
  D_{\cN}=\overline{D_{0}}
  \;\longrightarrow\;
  \delta_{\cN}=D_{\cN}^{*}.
\]
Given a matrix-positive, weighted-matrix-positive, form-closable
$\Gamma^{(2)}$-seed on a unital test algebra $\cA$ dense in
$L^{2}$, that paper produces a closed operator derivative and a
closed primitive divergence intrinsically; the construction is
universal in the inputs but conditional on their existence.

The conditional input is exactly what the present paper supplies,
in the simplest concrete nonclassical regime in which a
carr\'e-du-champ seed has plausible structure: the
\emph{diagonal} of the two-point motion of the Le
Jan--Raimond Brownian sticky flow~\cite{LeJanRaimond2004}.

\subsection{What is and is not claimed}\label{sec:intro-claims}

\textbf{Claimed.} Under the Le Jan--Raimond two-point
generator/domain convention (\cref{conv:LJR}), the diagonal
candidate
\begin{equation}\label{eq:Gamma-theta-intro}
  \Gamma^{(2)}_{\theta}(f,g)
  \;=\;\tfrac{c_{\theta}}{2}\,f'(x)g'(x)\,\mu_{\Delta}(dx),
  \qquad f,g\in\cT^{(2)},
\end{equation}
satisfies the five axioms required
by~\cite{GammaHierarchyDivergence}:
\textsc{(Sym)},
\textsc{(Vanish)},
\textsc{(MatPos)},
\textsc{(WMatPos)}, and
\textsc{(FormClos)} (\cref{thm:axioms}). By the transfer theorem
of~\cite{GammaHierarchyDivergence}, this yields a closed pair
$(\Dth,\divth)$ on $L^{2}(\R,\mu_{\Delta})$ with
integration-by-parts (\cref{thm:transfer}).

\textbf{Not claimed.}
\begin{itemize}
\item \emph{No flow-noise lift.} The pair $(\Dth,\divth)$ acts on
the diagonal Hilbert space $L^{2}(\R,\mu_{\Delta})$, not on
cylindrical functionals of the full flow noise $\cN_{\theta}$.
The Tier $\beta$ lift to flow-noise functionals is the next
input-existence problem and is out of scope (\cref{sec:tier-beta}).
\item \emph{No coefficient extraction.} The diagonal coefficient
$c_{\theta}$, the parametrization of $\theta$, the choice of
diagonal reference measure $\mu_{\Delta}$, and the precise
Wentzell-type boundary condition characterizing $L^{(2)}_{\theta}$
are all imported from the LJR convention. This paper does not
re-derive them, and it does not assert numerical values; we work
with $c_{\theta}\ge 0$ as a parameter tied to \cref{conv:LJR}.
The Tier $\alpha$ result is structural and survives any
self-consistent choice within the LJR family.
\item \emph{No black-noise universality.} This paper concerns the
LJR sticky/coalescing setting only. It says nothing about
Brownian web noise~\cite{SSS2017,Tsirelson2004}, two-dimensional
critical percolation noise~\cite{SchrammSmirnov2011}, or any
other black noise.
\end{itemize}

\subsection{Tier \texorpdfstring{$\alpha$}{alpha} versus Tier \texorpdfstring{$\beta$}{beta}}\label{sec:intro-tiers}

The companion construction~\cite{GammaHierarchyDivergence}
applies whenever a matrix-positive, weighted-matrix-positive,
form-closable $\Gamma^{(2)}$ on a unital test algebra dense in
$L^{2}(\Omega)$ is supplied. The probability space $\Omega$ and
the test algebra $\cA$ are inputs; the Tier ($\alpha/\beta$)
distinction reflects how rich an $\Omega$ one targets:
\begin{itemize}
\item[\textbf{$\alpha$.}] $\Omega=\R$, $\mu_{\Delta}$ as
reference measure, $\cA=\cT^{(2)}$ a smooth test class on the
diagonal of $\R^{2}$ identified with $\R$. The output is an
operator-derivative pair $(\Dth,\divth)$ on
$L^{2}(\R,\mu_{\Delta})$. This is the clean, fully provable
case, and is the content of the present paper.
\item[\textbf{$\beta$.}] $\Omega=$ the LJR flow noise
$\cN_{\theta}$, $\cA=$ a cylindrical algebra of functionals of
the flow $\Phi_{\theta}$. The output would be $(D,\delta)$ on
$L^{2}(\cN_{\theta})$. This requires a separate lift step that
we do not undertake here.
\end{itemize}

\subsection{Roadmap}\label{sec:intro-roadmap}

\Cref{sec:LJR-convention} fixes the LJR two-point
generator/domain convention as a parametrized statement
(\cref{conv:LJR}) and lists exactly which LJR data this paper
imports. \Cref{sec:product-defect} computes the two-point
product defect $\mathfrak D^{(2)}_{\theta}$ off the diagonal
(where it vanishes) and on the diagonal (where it is
$c_{\theta}f'(x)g'(x)\,\mu_{\Delta}(dx)$); the diagonal
candidate $\Gamma^{(2)}_{\theta}=\frac{1}{2}\mathfrak D^{(2)}_{\theta}$
is then defined. \Cref{sec:axioms} verifies the five
\textsc{GammaHierarchy} axioms. \Cref{sec:transfer} invokes the
transfer theorem of~\cite{GammaHierarchyDivergence} to obtain
the closed pair $(\Dth,\divth)$ and the integration-by-parts
identity. \Cref{sec:tier-beta} sketches the Tier $\beta$
problem and lists the open subproblems it presents.

%==============================================================
\section{LJR two-point sticky generator: convention}\label{sec:LJR-convention}
%==============================================================

\subsection{The one-point generator}\label{sec:LJR-onepoint}

We work on $(\R,\mathrm{Borel}(\R))$ throughout. The one-point
test class is
\[
  \cT^{(1)}\;:=\;C_{c}^{\infty}(\R),
\]
and the one-point generator is the standard half-Laplacian
\begin{equation}\label{eq:L1}
  L^{(1)}\;=\;\tfrac{1}{2}\partial_{x}^{2}\;\;\colon\;\;\cT^{(1)}\to C_{b}(\R).
\end{equation}
This is the marginal of the Brownian sticky flow $\Phi_{\theta}$
in every standard LJR parametrization; the assertion is
uncontroversial and is taken as given.

\subsection{Off-diagonal two-point generator}\label{sec:LJR-offdiag}

The two-point test class is a subspace
$\cT^{(2)}\subset C_{c}^{\infty}(\R^{2})$ to be specified by the
Wentzell condition below, and chosen so that the product
observables $f\otimes g$ with $f,g\in\cT^{(1)}$ lie in
$\cT^{(2)}$. Off the diagonal $\Delta:=\{(x,y)\in\R^{2}:x=y\}$,
the LJR two-point sticky flow evolves as two independent
Brownian motions:
\begin{equation}\label{eq:L2-offdiag}
  L^{(2)}_{\theta}\,F(x,y)\;=\;\tfrac{1}{2}\bigl(\partial_{x}^{2}+\partial_{y}^{2}\bigr)F(x,y),
  \qquad (x,y)\in\R^{2}\setminus\Delta,\;F\in\cT^{(2)}.
\end{equation}
This is the off-diagonal-independence property of LJR sticky
flows, valid in every parametrization of the family.

\subsection{Diagonal Wentzell condition: the LJR convention}\label{sec:LJR-diag}

The diagonal interaction is parametrized; we extract from LJR
exactly four ingredients and adopt them as a convention.

\begin{convention}[LJR two-point sticky generator/domain]\label{conv:LJR}
We adopt the following input data, all inherited from
\cite{LeJanRaimond2004}.
\begin{itemize}
\item[$(\mathrm a)$] \emph{Sticky parameter.}
$\theta\in[0,\infty]$ is the LJR stickiness parameter, with one
endpoint identifying the flow of independent Brownian motions
and the other the coalescing Brownian flow of
Arratia~\cite{LeJanRaimond2004}. We do not fix here which
endpoint is which; the discussion below applies symmetrically
once the LJR endpoint identification is fixed.
\item[$(\mathrm b)$] \emph{Diagonal coefficient.} A coefficient
$c_{\theta}\ge 0$, determined by the LJR Wentzell-type boundary
condition characterizing $L^{(2)}_{\theta}$ on
$\Dom(L^{(2)}_{\theta})\subset\cT^{(2)}$, controls the diagonal
interaction. The two endpoints of the parametrization satisfy
$c_{\theta_{\mathrm{indep}}}=0$ (independent flow) and
$c_{\theta_{\mathrm{coal}}}=c_{\mathrm{coal}}>0$ (coalescing
flow); intermediate $\theta$ give intermediate $c_{\theta}$.
\item[$(\mathrm c)$] \emph{Diagonal reference measure.} A Radon
measure $\mu_{\Delta}$ on $\Delta$, equivalently a Radon measure
on $\R$ via the identification $\Delta\simeq\R$, in which the
diagonal contribution to the two-point product defect is
expressed. Standard choices in the LJR literature are
Lebesgue measure on $\R$, the local-time-at-$0$ measure of the
difference $Y_{t}-X_{t}$ pushed forward to the diagonal, or a
multiple of one of these; the exact LJR choice and any rescaling
is part of \cref{conv:LJR}, and the final $c_{\theta}$ tracks
the choice.
\item[$(\mathrm d)$] \emph{Wentzell-type boundary condition.}
$L^{(2)}_{\theta}$ is the closed self-adjoint operator on
$L^{2}(\R^{2},dx\,dy)$ with off-diagonal action~\eqref{eq:L2-offdiag}
and a Wentzell-type boundary condition on $\Delta$ characterized
by $\theta$. The condition is symmetric in $(x,y)$, and yields the
closed symmetric Markov form on $L^{2}(\R^{2},dx\,dy)$ whose
generator is $L^{(2)}_{\theta}$. We refer
to~\cite{LeJanRaimond2004} for the exact statement.
\end{itemize}
\end{convention}

\begin{remark}[What \cref{conv:LJR} delegates]\label{rem:conv-delegates}
\Cref{conv:LJR} delegates the following to LJR: the explicit
form of the Wentzell condition; the explicit value of
$c_{\theta}$; the parametrization choice (which endpoint is
which); the choice of $\mu_{\Delta}$; and the closability of the
two-point Dirichlet form on $L^{2}(\R^{2},dx\,dy)$ under that
convention. The Tier $\alpha$ verification of
\cref{sec:axioms} is structural and depends only on
$c_{\theta}\ge 0$ and on $\mu_{\Delta}$ being a Radon measure on
$\R$, together with the form-closability of the diagonal energy
$\frac{c_{\theta}}{2}\int(f')^{2}\,d\mu_{\Delta}$ on a smooth
test class. None of these structural ingredients depends on the
numerical value of $c_{\theta}$ or on the precise Wentzell
condition.
\end{remark}

\subsection{Test class}\label{sec:LJR-testclass}

We denote by
\[
  \cT_{\Delta}\;:=\;C_{c}^{\infty}(\R)
\]
the smooth test class on $\R$, identified with the diagonal
$\Delta$ via $(x,x)\leftrightarrow x$. The diagonal carr\'e-du-champ
seed $\Gamma^{(2)}_{\theta}$ of \cref{sec:product-defect} will
be defined as a bilinear form
$\Gamma^{(2)}_{\theta}\colon\cT_{\Delta}\times\cT_{\Delta}\to L^{1}(\R,\mu_{\Delta})$,
and the test algebra fed into~\cite{GammaHierarchyDivergence} is
\begin{equation}\label{eq:test-algebra}
  \cA_{\theta}\;:=\;C_{c}^{\infty}(\R)\;\subset\;L^{\infty}(\R,\mu_{\Delta})\cap L^{2}(\R,\mu_{\Delta}),
\end{equation}
which is unital (after adjoining the constants in
$L^{\infty}(\R,\mu_{\Delta})$, see \cref{rem:unital}) and dense
in $L^{2}(\R,\mu_{\Delta})$ for any Radon $\mu_{\Delta}$ that is
$\sigma$-finite and supported on $\R$.

\begin{remark}[Unitality of $\cA_{\theta}$]\label{rem:unital}
$C_{c}^{\infty}(\R)$ does not contain the constant function $1$
when $\mu_{\Delta}(\R)=\infty$. We work with the unital
extension $\cA_{\theta}^{u}:=\R\cdot 1+C_{c}^{\infty}(\R)$ when
the unitality hypothesis of \cite[\S 2]{GammaHierarchyDivergence}
is needed; the carr\'e-du-champ
$\Gamma^{(2)}_{\theta}(c+f,c'+g)=\Gamma^{(2)}_{\theta}(f,g)$ by
\textsc{(Vanish)} below, so the extension is consistent. To
lighten notation we write $\cA_{\theta}$ for the unital extension
when no confusion arises.
\end{remark}

%==============================================================
\section{Product defect and diagonal carr\'e-du-champ seed}\label{sec:product-defect}
%==============================================================

\subsection{The two-point product defect}\label{sec:Dgamma}

For $f,g\in\cT^{(1)}=C_{c}^{\infty}(\R)$, the two-point product
defect of the LJR sticky flow is, in the notation of
\cite[Def.~6.7]{PrimitiveDivergence},
\begin{equation}\label{eq:Dgamma-def}
  \mathfrak D^{(2)}_{\theta}(f,g)
  \;:=\;
  L^{(2)}_{\theta}(f\otimes g)\;-\;(L^{(1)}f)\otimes g\;-\;f\otimes(L^{(1)}g),
\end{equation}
read as a distribution on $\R^{2}$ obtained from the closed
generator $L^{(2)}_{\theta}$ of \cref{conv:LJR}.

\begin{lemma}[Off-diagonal vanishing]\label{lem:offdiag}
For $f,g\in C_{c}^{\infty}(\R)$ and $(x,y)\in\R^{2}\setminus\Delta$,
\begin{equation}\label{eq:offdiag-zero}
  \mathfrak D^{(2)}_{\theta}(f,g)(x,y)\;=\;0.
\end{equation}
\end{lemma}

\begin{proof}
By the off-diagonal independence~\eqref{eq:L2-offdiag} of
\cref{conv:LJR}$(\mathrm a)$,
\[
  L^{(2)}_{\theta}(f\otimes g)(x,y)
  \;=\;\tfrac{1}{2}f''(x)g(y)+\tfrac{1}{2}f(x)g''(y),
  \qquad x\neq y,
\]
which equals $(L^{(1)}f)(x)g(y)+f(x)(L^{(1)}g)(y)$ by
\eqref{eq:L1}. Subtracting gives~\eqref{eq:offdiag-zero}.
\end{proof}

\subsection{Diagonal contribution}\label{sec:diag-contribution}

The diagonal contribution to $\mathfrak D^{(2)}_{\theta}$ is the
content of the Wentzell-type condition of
\cref{conv:LJR}$(\mathrm d)$. Reading the LJR convention as a
parametric statement, we record:

\begin{proposition}[Diagonal contribution under \cref{conv:LJR}]\label{prop:diag-contribution}
Under \cref{conv:LJR}, the two-point product defect
$\mathfrak D^{(2)}_{\theta}(f,g)$ is a distribution on $\R^{2}$
supported on the diagonal $\Delta$, with diagonal density
\begin{equation}\label{eq:Dgamma-diag}
  \mathfrak D^{(2)}_{\theta}(f,g)
  \;=\;c_{\theta}\,f'(x)g'(x)\,\mu_{\Delta}(dx)
  \quad\text{on }\Delta,
\end{equation}
where $c_{\theta}\ge 0$ is the diagonal interaction coefficient
and $\mu_{\Delta}$ is the diagonal reference measure of
\cref{conv:LJR}$(\mathrm b),(\mathrm c)$. Equivalently, for
$f,g,\varphi\in C_{c}^{\infty}(\R^{2})$ with
$\varphi$ a test function,
\begin{equation}\label{eq:Dgamma-pairing}
  \langle\mathfrak D^{(2)}_{\theta}(f,g),\varphi\rangle
  \;=\;c_{\theta}\!\int_{\R}\!f'(x)g'(x)\,\varphi(x,x)\,\mu_{\Delta}(dx).
\end{equation}
\end{proposition}

\begin{proof}
\Cref{lem:offdiag} shows that $\mathfrak D^{(2)}_{\theta}(f,g)$,
viewed as a distribution on $\R^{2}$, is supported on $\Delta$
when read against test functions on $\R^{2}$. The remaining
diagonal pairing~\eqref{eq:Dgamma-pairing} is the content of
\cref{conv:LJR}: by the Wentzell-type characterization of
$L^{(2)}_{\theta}$, the diagonal contribution is bilinear in
$(f,g)$, vanishes when either $f$ or $g$ is constant
(consistent with the marginals~\eqref{eq:L1}), and is
quadratic-in-derivative on the diagonal because the off-diagonal
generator is the half-Laplacian. The two-point sticky generator
under \cref{conv:LJR} prescribes this pairing as
\eqref{eq:Dgamma-pairing} with coefficient $c_{\theta}\ge 0$ and
reference measure $\mu_{\Delta}$. The structural form is uniform
across the LJR parametrization; the specific
$(c_{\theta},\mu_{\Delta})$ pair is what
\cref{conv:LJR}$(\mathrm b),(\mathrm c)$ delegate.
\end{proof}

\begin{remark}[Coalescing endpoint cross-check]\label{rem:coalescing-check}
In the coalescing endpoint of the LJR family, the two-point
generator on the diagonal is
$L^{(2)}(f\otimes g)(x,x)=L^{(1)}(fg)(x)=\tfrac{1}{2}(fg)''(x)$,
and a direct computation
(\cite[Prop.~6.10]{PrimitiveDivergence})
gives
$\mathfrak D^{(2)}(f,g)(x,x)=f'(x)g'(x)$,
i.e.\ \eqref{eq:Dgamma-diag} with
$c_{\mathrm{coal}}=1$ and $\mu_{\Delta}=\delta_{\Delta}$ (the
diagonal Dirac), this last being the standard cross-section
identification with diagonal Lebesgue measure. The exact
$c_{\mathrm{coal}}$ in the LJR sticky-family parametrization is
inherited from \cref{conv:LJR}$(\mathrm a),(\mathrm c)$.
\end{remark}

\subsection{The diagonal carr\'e-du-champ seed}\label{sec:Gamma-seed}

\begin{definition}[Diagonal carr\'e-du-champ seed]\label{def:Gamma-theta}
The \emph{diagonal carr\'e-du-champ seed} associated with the
LJR sticky flow $\Phi_{\theta}$ under \cref{conv:LJR} is
\begin{equation}\label{eq:Gamma-theta}
  \Gamma^{(2)}_{\theta}(f,g)
  \;:=\;\tfrac{1}{2}\,\mathfrak D^{(2)}_{\theta}(f,g)
  \;=\;\frac{c_{\theta}}{2}\,f'(x)g'(x)\,\mu_{\Delta}(dx)
  \quad\text{on }\Delta,
\end{equation}
read as a measure-valued bilinear form on
$\cT_{\Delta}\times\cT_{\Delta}$.
\end{definition}

The factor $\frac{1}{2}$ is the Bakry--\'Emery normalization
\cite[Eq.~(5.7)]{PrimitiveDivergence} of the
$\Gamma$-functional from the product defect, and matches the
classical formula $\Gamma(F,G)=\frac{1}{2}(L(FG)-FLG-GLF)$ at
the level of the one-point semigroup
(\cite[Ch.~1]{BakryGentilLedoux}).

\begin{remark}[Reading as $L^{1}(\R,\mu_{\Delta})$]\label{rem:Gamma-as-L1}
For $f,g\in\cT_{\Delta}=C_{c}^{\infty}(\R)$, the right-hand side
of~\eqref{eq:Gamma-theta} is the Radon measure
$\frac{c_{\theta}}{2}f'g'\,d\mu_{\Delta}$ on $\R$. Since
$f',g'\in L^{\infty}(\R)$ are compactly supported, the density
$\frac{c_{\theta}}{2}f'(x)g'(x)$ is in
$L^{\infty}(\R,\mu_{\Delta})\cap L^{1}(\R,\mu_{\Delta})$. We
will accordingly read $\Gamma^{(2)}_{\theta}(f,g)$ as an element
of $L^{1}(\R,\mu_{\Delta})$, with the implicit identification
of the measure with its Radon--Nikod\'ym density against
$\mu_{\Delta}$:
\begin{equation}\label{eq:Gamma-L1}
  \Gamma^{(2)}_{\theta}(f,g)\;=\;\tfrac{c_{\theta}}{2}\,f'(x)g'(x)
  \quad\text{in }L^{1}(\R,\mu_{\Delta}).
\end{equation}
This is the form in which $\Gamma^{(2)}_{\theta}$ is fed into
the construction of \cite{GammaHierarchyDivergence}, where the
ambient probability space is $(\R,\mu_{\Delta}/\mu_{\Delta}(\R))$
when $\mu_{\Delta}$ is finite, and $(\R,\mu_{\Delta})$ as a
$\sigma$-finite measure space otherwise (the construction of
\cite{GammaHierarchyDivergence} extends to $\sigma$-finite
measure spaces with no change in the proofs, since matrix and
weighted-matrix positivity are pointwise-on-$\R$ statements).
\end{remark}

%==============================================================
\section{Verification of the GammaHierarchy axioms}\label{sec:axioms}
%==============================================================

This section verifies the five axioms of
\cite[Ass.~2.4]{GammaHierarchyDivergence} (under the renaming
\textsc{Sym, Vanish, MatPos, WMatPos, FormClos} of the
matrix-positivity-and-form-closability axiom set) for the
diagonal seed $\Gamma^{(2)}_{\theta}$ of
\cref{def:Gamma-theta}, on the test algebra $\cA_{\theta}$
of~\eqref{eq:test-algebra}.

\subsection{Symmetry and constancy vanishing}\label{sec:axioms-symvan}

\begin{lemma}[\normalfont\textsc{(Sym)} and \normalfont\textsc{(Vanish)}]\label{lem:symvan}
For $f,g\in\cA_{\theta}$ and $c\in\R$:
\begin{align}
  \Gamma^{(2)}_{\theta}(f,g) &\;=\;\Gamma^{(2)}_{\theta}(g,f)
  \quad\text{in }L^{1}(\R,\mu_{\Delta}),\label{eq:axsym}\\
  \Gamma^{(2)}_{\theta}(c,g) &\;=\;0
  \quad\text{in }L^{1}(\R,\mu_{\Delta}).\label{eq:axvan}
\end{align}
\end{lemma}

\begin{proof}
By~\eqref{eq:Gamma-L1},
$\Gamma^{(2)}_{\theta}(f,g)=\frac{c_{\theta}}{2}f'(x)g'(x)
=\frac{c_{\theta}}{2}g'(x)f'(x)=\Gamma^{(2)}_{\theta}(g,f)$,
which is~\eqref{eq:axsym}. For~\eqref{eq:axvan}: $c'=0$ for any
constant $c\in\R$, so
$\Gamma^{(2)}_{\theta}(c,g)=\frac{c_{\theta}}{2}\cdot 0\cdot g'(x)=0$
in $L^{1}(\R,\mu_{\Delta})$.
\end{proof}

\subsection{Matrix positivity}\label{sec:axioms-matpos}

\begin{lemma}[\normalfont\textsc{(MatPos)}]\label{lem:matpos}
For every finite collection $f_{1},\ldots,f_{n}\in\cA_{\theta}$
and $U_{1},\ldots,U_{n}\in\cA_{\theta}$,
\begin{equation}\label{eq:axmatpos}
  \sum_{i,j=1}^{n}\!\int_{\R}U_{i}(x)U_{j}(x)\,\Gamma^{(2)}_{\theta}(f_{i},f_{j})(x)\,\mu_{\Delta}(dx)\;\ge\;0.
\end{equation}
Equivalently, in the notation of
\cite[Eq.~(2.3)]{GammaHierarchyDivergence},
$\sum_{i,j}\E[U_{i}U_{j}\,\Gamma(f_{i},f_{j})]\ge 0$ on
$(\R,\mu_{\Delta})$.
\end{lemma}

\begin{proof}
By~\eqref{eq:Gamma-L1},
\[
  \sum_{i,j}U_{i}(x)U_{j}(x)\,\Gamma^{(2)}_{\theta}(f_{i},f_{j})(x)
  \;=\;\frac{c_{\theta}}{2}\sum_{i,j}U_{i}(x)U_{j}(x)\,f_{i}'(x)f_{j}'(x).
\]
The double sum factors as a square of a sum:
\begin{equation}\label{eq:square-of-sum}
  \sum_{i,j}U_{i}(x)U_{j}(x)\,f_{i}'(x)f_{j}'(x)
  \;=\;\Bigl(\sum_{i}U_{i}(x)\,f_{i}'(x)\Bigr)^{2}
  \;\ge\;0
\end{equation}
pointwise in $x\in\R$. Therefore
$\frac{c_{\theta}}{2}(\sum_{i}U_{i}f_{i}')^{2}\ge 0$ pointwise
on $\R$ (using $c_{\theta}\ge 0$ from \cref{conv:LJR}$(\mathrm b)$),
and integrating against the nonnegative Radon measure
$\mu_{\Delta}$ preserves the inequality.
\end{proof}

\begin{remark}[Pointwise PSD-kernel structure]\label{rem:pointwise-PSD}
\Cref{lem:matpos} is the textbook diffusion-carr\'e-du-champ
case described as automatic in
\cite[Rem.~2.6]{GammaHierarchyDivergence}: the kernel
$\Gamma^{(2)}_{\theta}(f,g)(x)=\frac{c_{\theta}}{2}f'(x)g'(x)$ is
pointwise positive-semidefinite on $\R$ (the
$f\otimes g\mapsto f'g'$ map factors through the gradient
$\nabla\colon\cA_{\theta}\to L^{2}(\R,\mu_{\Delta})$ given by
$\nabla f=f'$). This pointwise PSD-kernel structure
\emph{automatically} gives both \textsc{(MatPos)} and
\textsc{(WMatPos)}; see \cref{lem:wmatpos} below.
\end{remark}

\subsection{Weighted matrix positivity}\label{sec:axioms-wmatpos}

\begin{lemma}[\normalfont\textsc{(WMatPos)}]\label{lem:wmatpos}
For every finite collection $f_{1},\ldots,f_{n}\in\cA_{\theta}$,
every $U_{1},\ldots,U_{n}\in\cA_{\theta}$, and every nonnegative
multiplier $Q\in\cA_{\theta,+}$ (the cone of nonnegative elements
of $\cA_{\theta}$),
\begin{equation}\label{eq:axwmatpos}
  \sum_{i,j=1}^{n}\!\int_{\R}Q(x)\,U_{i}(x)U_{j}(x)\,\Gamma^{(2)}_{\theta}(f_{i},f_{j})(x)\,\mu_{\Delta}(dx)\;\ge\;0.
\end{equation}
\end{lemma}

\begin{proof}
By~\eqref{eq:Gamma-L1} and the square-of-sum
identity~\eqref{eq:square-of-sum},
\[
  Q(x)\,\sum_{i,j}U_{i}(x)U_{j}(x)\,\Gamma^{(2)}_{\theta}(f_{i},f_{j})(x)
  \;=\;\frac{c_{\theta}}{2}\,Q(x)\,\Bigl(\sum_{i}U_{i}(x)\,f_{i}'(x)\Bigr)^{2}.
\]
Since $Q(x)\ge 0$, $c_{\theta}\ge 0$, and the parenthesized
expression is squared, the integrand is pointwise nonnegative on
$\R$. Integrating against the nonnegative Radon measure
$\mu_{\Delta}$ preserves the inequality.
\end{proof}

\subsection{Form closability}\label{sec:axioms-formclos}

The remaining axiom is form-closability of the energy form
\begin{equation}\label{eq:cE-theta}
  \cE_{\theta}(f,g)\;:=\;\int_{\R}\Gamma^{(2)}_{\theta}(f,g)(x)\,\mu_{\Delta}(dx)
  \;=\;\frac{c_{\theta}}{2}\!\int_{\R}f'(x)g'(x)\,\mu_{\Delta}(dx),
  \qquad f,g\in\cA_{\theta},
\end{equation}
in the sense of \cite[\textsc{(FormClos)} of
Ass.~2.4]{GammaHierarchyDivergence}: if $(f_{n})\subset\cA_{\theta}$
satisfies $f_{n}\to 0$ in $L^{2}(\R,\mu_{\Delta})$ and
$\cE_{\theta}(f_{n}-f_{m},f_{n}-f_{m})\to 0$ as $n,m\to\infty$,
then $\cE_{\theta}(f_{n},f_{n})\to 0$.

\begin{lemma}[\normalfont\textsc{(FormClos)}]\label{lem:formclos}
Suppose, in addition to \cref{conv:LJR}, that the diagonal
reference measure $\mu_{\Delta}$ is a Radon measure on $\R$
under which the gradient
$\nabla\colon C_{c}^{\infty}(\R)\to L^{2}(\R,\mu_{\Delta})$,
$\nabla f=f'$, is closable in the unbounded-operator sense from
$L^{2}(\R,\mu_{\Delta})$ to $L^{2}(\R,\mu_{\Delta})$. Then the
energy form $\cE_{\theta}$ of~\eqref{eq:cE-theta} is
form-closable on $L^{2}(\R,\mu_{\Delta})$.
\end{lemma}

\begin{proof}
Suppose $(f_{n})\subset C_{c}^{\infty}(\R)$ satisfies
$f_{n}\to 0$ in $L^{2}(\R,\mu_{\Delta})$ and
$\cE_{\theta}(f_{n}-f_{m},f_{n}-f_{m})\to 0$. By
\eqref{eq:cE-theta},
\[
  \cE_{\theta}(f_{n}-f_{m},f_{n}-f_{m})
  \;=\;\frac{c_{\theta}}{2}\!\int_{\R}(f_{n}'-f_{m}')^{2}\,d\mu_{\Delta}
  \;=\;\frac{c_{\theta}}{2}\,\|f_{n}'-f_{m}'\|_{L^{2}(\mu_{\Delta})}^{2}.
\]
If $c_{\theta}=0$ there is nothing to prove ($\cE_{\theta}\equiv 0$).
If $c_{\theta}>0$, the hypothesis says
$\|f_{n}'-f_{m}'\|_{L^{2}(\mu_{\Delta})}\to 0$, so $(f_{n}')$ is
Cauchy in $L^{2}(\R,\mu_{\Delta})$ and converges to some
$g\in L^{2}(\R,\mu_{\Delta})$. The closability of the gradient
$\nabla$ from $L^{2}(\R,\mu_{\Delta})$ to itself, applied to the
sequence $(f_{n})$ with $f_{n}\to 0$ in $L^{2}(\R,\mu_{\Delta})$
and $\nabla f_{n}=f_{n}'\to g$ in $L^{2}(\R,\mu_{\Delta})$,
forces $g=0$. Hence
$\cE_{\theta}(f_{n},f_{n})=\frac{c_{\theta}}{2}\|f_{n}'\|_{L^{2}(\mu_{\Delta})}^{2}\to\frac{c_{\theta}}{2}\|g\|_{L^{2}(\mu_{\Delta})}^{2}=0$.
\end{proof}

\begin{remark}[When the gradient closability hypothesis holds]\label{rem:gradient-closability}
The closability hypothesis of \cref{lem:formclos} on the
gradient $\nabla\colon C_{c}^{\infty}(\R)\to L^{2}(\R,\mu_{\Delta})$
is the standard hypothesis on the diagonal measure
$\mu_{\Delta}$ from the theory of weighted Dirichlet forms
(\cite[Ch.~1]{FukushimaOshimaTakeda}). It holds in the standard
LJR cases: when $\mu_{\Delta}$ is Lebesgue measure on $\R$, the
gradient is the standard Sobolev gradient and is closable on
$L^{2}(\R,dx)$ with closure equal to $W^{1,2}(\R)\to L^{2}(\R,dx)$
(\cite[Ch.~1, \S 1.4]{FukushimaOshimaTakeda},
\cite[Ch.~1]{BakryGentilLedoux}); when $\mu_{\Delta}$ is a
local-time-type measure, the gradient closability follows from
the symmetric-Markov-generator framework of LJR
(\cite{LeJanRaimond2004}) together with the standard
correspondence between symmetric Markov semigroups and closable
Dirichlet forms. In each case, closability is part of the
LJR convention package; we mark it explicitly in
\cref{lem:formclos} because it is the only axiom that is not
purely algebraic and pointwise.
\end{remark}

\subsection{Main verification theorem}\label{sec:axioms-main}

\begin{theorem}[The diagonal seed satisfies the GammaHierarchy axioms]\label{thm:axioms}
Under \cref{conv:LJR} together with the gradient-closability
hypothesis of \cref{lem:formclos}, the diagonal seed
$\Gamma^{(2)}_{\theta}\colon\cA_{\theta}\times\cA_{\theta}\to L^{1}(\R,\mu_{\Delta})$
of \cref{def:Gamma-theta} satisfies the five axioms required by
\cite[Ass.~2.4]{GammaHierarchyDivergence}:
{\normalfont\textsc{(Sym)}} (\cref{lem:symvan}), {\normalfont\textsc{(Vanish)}}
(\cref{lem:symvan}), {\normalfont\textsc{(MatPos)}} (\cref{lem:matpos}),
{\normalfont\textsc{(WMatPos)}} (\cref{lem:wmatpos}), and {\normalfont\textsc{(FormClos)}}
(\cref{lem:formclos}).
\end{theorem}

\begin{proof}
Each axiom is the corresponding lemma of this section.
\end{proof}

%==============================================================
\section{Construction by transfer}\label{sec:transfer}
%==============================================================

\Cref{thm:axioms} verifies the input hypotheses required by
\cite[Thm.~A and Thm.~B]{GammaHierarchyDivergence} on the test
algebra $\cA_{\theta}$ over the measure space
$(\R,\mu_{\Delta})$. The transfer is immediate.

\begin{theorem}[Transfer to the GammaHierarchy construction]\label{thm:transfer}
Under \cref{conv:LJR} and the gradient-closability hypothesis of
\cref{lem:formclos}, applying
\cite[Thm.~A and Thm.~B]{GammaHierarchyDivergence} to the
diagonal seed $\Gamma^{(2)}_{\theta}$ of \cref{def:Gamma-theta}
on the test algebra $\cA_{\theta}\subset L^{\infty}(\R,\mu_{\Delta})\cap L^{2}(\R,\mu_{\Delta})$
yields:
\begin{enumerate}
\item[$(\mathrm i)$] a Hilbert one-form module
$\cU_{\Gamma_{\theta}}=\overline{\cA_{\theta}\otimes\cA_{\theta}/\mathcal N_{0}}$,
with semi-inner product
\[
  \langle [F\otimes U],[G\otimes V]\rangle_{\cU_{\Gamma_{\theta}}}
  \;=\;\frac{c_{\theta}}{2}\!\int_{\R}\!U(x)V(x)\,F'(x)G'(x)\,\mu_{\Delta}(dx),
\]
together with the bounded $\cA_{\theta}$-bimodule action
$W\cdot[F\otimes U]:=[F\otimes(WU)]$ for $W\in\cA_{\theta}$,
satisfying $\|W\cdot u\|\le\|W\|_{\infty}\|u\|$;
\item[$(\mathrm{ii})$] a closable seed
$D_{0}^{\theta}\colon\cA_{\theta}\to\cU_{\Gamma_{\theta}}$,
$D_{0}^{\theta}F:=[F\otimes 1]$, with isometry
$\langle D_{0}^{\theta}F,D_{0}^{\theta}G\rangle_{\cU_{\Gamma_{\theta}}}
=\frac{c_{\theta}}{2}\int_{\R}F'G'\,d\mu_{\Delta}=\cE_{\theta}(F,G)$;
\item[$(\mathrm{iii})$] a closed densely-defined operator
derivative
\[
  \Dth\;:=\;\overline{D_{0}^{\theta}}\;\;\colon\;\;\Dom(\Dth)\subset L^{2}(\R,\mu_{\Delta})\to\cU_{\Gamma_{\theta}};
\]
\item[$(\mathrm{iv})$] a closed densely-defined primitive
divergence
\[
  \divth\;:=\;\Dth^{*}\;\;\colon\;\;\Dom(\divth)\subset\cU_{\Gamma_{\theta}}\to L^{2}(\R,\mu_{\Delta}),
\]
with the explicit adjoint domain
\[
  \Dom(\divth)
  =
  \left\{
  u\in\cU_{\Gamma_{\theta}}:
  F\mapsto\langle\Dth F,u\rangle_{\cU_{\Gamma_{\theta}}}
  \text{ is }L^{2}(\mu_{\Delta})\text{-continuous on }\Dom(\Dth)
  \right\}.
\]
\item[$(\mathrm v)$] the integration-by-parts identity
\begin{equation}\label{eq:IBP-theta}
  \int_{\R}F\,\divth(u)\,d\mu_{\Delta}
  \;=\;\langle\Dth F,u\rangle_{\cU_{\Gamma_{\theta}}},
  \qquad F\in\Dom(\Dth),\;u\in\Dom(\divth).
\end{equation}
\end{enumerate}
The pair $(\Dth,\divth)$ is the operator-derivative pair on
$L^{2}(\R,\mu_{\Delta})$ associated with the diagonal LJR
sticky-flow seed.
\end{theorem}

\begin{proof}
\Cref{thm:axioms} verifies \cref{conv:LJR}'s data fits the
hypotheses of \cite[Ass.~2.4]{GammaHierarchyDivergence}.
\Cref{thm:transfer}$(\mathrm i),(\mathrm{ii})$ are then
\cite[Thm.~A]{GammaHierarchyDivergence} read on $\cA_{\theta}$;
$(\mathrm{iii}),(\mathrm{iv}),(\mathrm v)$ are
\cite[Thm.~B]{GammaHierarchyDivergence}.
\end{proof}

\begin{remark}[Trivial endpoint at $c_{\theta}=0$]\label{rem:trivial-endpoint}
At the LJR independent-flow endpoint, $c_{\theta}=0$ by
\cref{conv:LJR}$(\mathrm b)$. Then
$\Gamma^{(2)}_{\theta}\equiv 0$, $\cE_{\theta}\equiv 0$,
$\cU_{\Gamma_{\theta}}=\{0\}$, and $\Dth\equiv 0$,
$\divth\equiv 0$. This is consistent with the
trivial-product-defect statement at the independent endpoint
(\cite[Prop.~6.5]{PrimitiveDivergence}): no operator derivative
content is generated by the diagonal because the diagonal
interaction is empty.
\end{remark}

\begin{remark}[Coalescing endpoint]\label{rem:coalescing-endpoint}
At the LJR coalescing endpoint, $c_{\theta}=c_{\mathrm{coal}}>0$.
The construction yields a nontrivial pair
$(\Dth,\divth)$ on $L^{2}(\R,\mu_{\Delta})$, with energy form
$\cE_{\theta}(f,f)=\frac{c_{\mathrm{coal}}}{2}\int(f')^{2}\,d\mu_{\Delta}$.
With the standard choice $\mu_{\Delta}=dx$ (Lebesgue), this is
$\cE_{\theta}(f,f)=\frac{c_{\mathrm{coal}}}{2}\|f'\|_{L^{2}(\R)}^{2}$,
i.e.\ a multiple of the standard Dirichlet form on $\R$, and
$\Dth$ is a multiple of the standard Sobolev gradient
$W^{1,2}(\R)\to L^{2}(\R)$, with $\divth$ the corresponding
distributional divergence.
\end{remark}

\begin{remark}[Generator $L_{\theta}=-\divth\Dth$]\label{rem:Ltheta}
Following \cite[Rem.~4.10]{GammaHierarchyDivergence}, the
composition
$L_{\theta}:=-\divth\Dth=-\Dth^{*}\Dth$ on
$\Dom(L_{\theta}):=\{F\in\Dom(\Dth):\Dth F\in\Dom(\divth)\}$
is self-adjoint and nonpositive on $L^{2}(\R,\mu_{\Delta})$, by
the standard $T^{*}T$ theorem
(\cite[Thm.~13.13]{Rudin-FA}). With the standard Lebesgue
choice $\mu_{\Delta}=dx$, $L_{\theta}=\frac{c_{\theta}}{2}\partial_{x}^{2}$
on the natural Sobolev domain $W^{2,2}(\R)$, and the associated
symmetric strongly-continuous contraction semigroup is
$\frac{c_{\theta}}{2}$-rescaled heat. With other choices of
$\mu_{\Delta}$ inherited from \cref{conv:LJR}, $L_{\theta}$ is
the corresponding weighted Laplacian on $L^{2}(\R,\mu_{\Delta})$.
\end{remark}

%==============================================================
\section{Tier \texorpdfstring{$\beta$}{beta} outlook}\label{sec:tier-beta}
%==============================================================

\subsection{What Tier \texorpdfstring{$\alpha$}{alpha} did}\label{sec:beta-recap}

\Cref{thm:transfer} produces the operator-derivative pair
$(\Dth,\divth)$ on the diagonal Hilbert space
$L^{2}(\R,\mu_{\Delta})$. The probability space underlying the
GammaHierarchy construction is $(\R,\mu_{\Delta})$ (after
normalization, when $\mu_{\Delta}$ is finite); the test algebra
is $\cA_{\theta}=C_{c}^{\infty}(\R)$. This is the
\emph{diagonal} of the two-point motion of the LJR sticky flow,
not the flow itself.

\subsection{What Tier \texorpdfstring{$\beta$}{beta} would do}\label{sec:beta-target}

A complete operator-derivative theory of the LJR sticky flow
would produce a closed pair $(D,\delta)$ on $L^{2}(\cN_{\theta})$,
where $\cN_{\theta}$ is the noise of the flow $\Phi_{\theta}$
in the Tsirelson--Vershik sense
(\cite{Tsirelson2004,LeJanRaimond2004},
\cite[Def.~2.4]{PrimitiveDivergence}). The natural test algebra
$\cA_{\cN_{\theta}}$ is a cylindrical algebra of bounded
functionals of the flow; standard choices include polynomial
combinations of evaluation functionals
$F(\Phi)=\langle\mu_{x}^{\Phi},\phi\rangle$ for finitely many
starting points $x_{i}\in\R$ and Schwartz $\phi_{i}$, where
$\mu_{x}^{\Phi}$ is the random measure produced by
$\Phi_{\theta}$ from the initial point $x$.

The Tier $\beta$ problem is to lift the diagonal seed
$\Gamma^{(2)}_{\theta}$ of \cref{def:Gamma-theta} to a
$\Gamma^{(2)}_{\cN_{\theta}}\colon\cA_{\cN_{\theta}}\times\cA_{\cN_{\theta}}\to L^{1}(\cN_{\theta})$
satisfying the five axioms of
\cite[Ass.~2.4]{GammaHierarchyDivergence}, and then to apply the
construction of \cite{GammaHierarchyDivergence} on
$L^{2}(\cN_{\theta})$.

\subsection{Subproblems for Tier \texorpdfstring{$\beta$}{beta}}\label{sec:beta-subproblems}

The lift step decomposes into the following subproblems, in
roughly increasing order of difficulty.

\begin{enumerate}
\item[$(\mathrm B_{1})$] \emph{Choice of cylindrical algebra.}
Specify a unital subalgebra
$\cA_{\cN_{\theta}}\subset L^{\infty}(\cN_{\theta})\cap L^{2}(\cN_{\theta})$
dense in $L^{2}(\cN_{\theta})$, closed under products, and
containing enough functionals of the LJR flow that the lift of
$\Gamma^{(2)}_{\theta}$ is well defined. The natural candidates
are polynomial-in-evaluations algebras as described above.
\item[$(\mathrm B_{2})$] \emph{Lift formula.} For
$F,G\in\cA_{\cN_{\theta}}$, define
$\Gamma^{(2)}_{\cN_{\theta}}(F,G)$ from the diagonal seed
$\Gamma^{(2)}_{\theta}$ by averaging the diagonal pairing
\eqref{eq:Dgamma-pairing} against the random measure-valued
data of $\Phi_{\theta}$. The natural candidate is, schematically,
\[
  \Gamma^{(2)}_{\cN_{\theta}}(F,G)
  \;\sim\;c_{\theta}\!\int_{\Phi_{\theta}\text{-paths}}\!\!\langle DF,DG\rangle\,d\mathbb{P}^{\Phi_{\theta}},
\]
where $DF$ is a $\cN_{\theta}$-Malliavin-type derivative
adapted to the LJR-sticky structure. Pinning this down on the
chosen $\cA_{\cN_{\theta}}$ is the central technical step.
\item[$(\mathrm B_{3})$] \emph{(MatPos), (WMatPos) at the
flow-noise level.} The pointwise PSD-kernel structure of
\cref{rem:pointwise-PSD} no longer applies automatically once
one passes from the diagonal of $\R^{2}$ to the flow-path
space. Whether the lifted $\Gamma^{(2)}_{\cN_{\theta}}$ inherits
matrix and weighted-matrix positivity from the diagonal seed
needs separate verification, and may require structural
hypotheses on the chosen cylindrical algebra.
\item[$(\mathrm B_{4})$] \emph{(FormClos) at the flow-noise
level.} Form-closability of the lifted energy form on
$L^{2}(\cN_{\theta})$ is the deepest of the four. It is closely
related to closability of the LJR generator on the
flow-noise level.
\item[$(\mathrm B_{5})$] \emph{Identification of the lifted
$(D,\delta)$.} Assuming $(\mathrm B_{1})$--$(\mathrm B_{4})$, the
closed pair $(D_{\cN_{\theta}},\delta_{\cN_{\theta}})$ produced
by \cite{GammaHierarchyDivergence} on $L^{2}(\cN_{\theta})$
should be identified with a Malliavin-type pair adapted to the
LJR-sticky structure, in the sense of
\cite[Ch.~6]{Nualart} for the classical case. This step
relates the abstract construction to existing Malliavin-style
calculi for stochastic flows.
\end{enumerate}

The present paper does not attempt $(\mathrm B_{1})$--$(\mathrm B_{5})$;
it produces only the diagonal-level $(\Dth,\divth)$ of
\cref{thm:transfer}. The diagonal seed and its axioms
established here are the input data the Tier $\beta$ lift would
operate on.

\subsection{Closing summary}\label{sec:beta-summary}

This paper supplies the first concrete input-existence step in
the operator-derivative program over a nonclassical-noise
candidate. Working under a stated LJR convention
(\cref{conv:LJR}), we extracted a diagonal carr\'e-du-champ seed
$\Gamma^{(2)}_{\theta}$ from the LJR sticky/coalescing flow and
verified the five axioms required by
\cite{GammaHierarchyDivergence}. The transfer
(\cref{thm:transfer}) yields a closed operator-derivative pair
$(\Dth,\divth)$ on $L^{2}(\R,\mu_{\Delta})$. The full
flow-noise lift $(D_{\cN_{\theta}},\delta_{\cN_{\theta}})$ is the
next input-existence problem.

%==============================================================
\begin{thebibliography}{99}

\bibitem{GammaHierarchyDivergence}
R.~Fontes,
\emph{From Carr\'e-du-Champ Hierarchies to Primitive Divergence}.
Preprint.

\bibitem{PrimitiveDivergence}
R.~Fontes,
\emph{Primitive Divergence and Carr\'e-du-Champ Seeds over Nonclassical Noises}.
Preprint.

\bibitem{EnergyFirst}
R.~Fontes,
\emph{Energy-First Stochastic Calculus: Iterated Brownian Motion as a Test Case}.
Preprint.

\bibitem{FCDC}
R.~Fontes,
\emph{First-Chaos Dual Closure of the IBM It\^o Correction}.
Preprint.

\bibitem{BidegreeAtlas}
R.~Fontes,
\emph{A Bidegree Atlas for the IBM It\^o Correction}.
Preprint.

\bibitem{BakryGentilLedoux}
D.~Bakry, I.~Gentil and M.~Ledoux,
\emph{Analysis and Geometry of Markov Diffusion Operators}.
Grundlehren der mathematischen Wissenschaften, vol.~348, Springer, 2014.

\bibitem{FukushimaOshimaTakeda}
M.~Fukushima, Y.~\=Oshima and M.~Takeda,
\emph{Dirichlet Forms and Symmetric Markov Processes}, 2nd revised ed.
de Gruyter Studies in Mathematics, vol.~19, de Gruyter, Berlin, 2011.

\bibitem{LeJanRaimond2004}
Y.~Le~Jan and O.~Raimond,
Flows, coalescence and noise.
\emph{Ann. Probab.} \textbf{32} (2004), no.~2A, 1247--1315.

\bibitem{Nualart}
D.~Nualart,
\emph{The Malliavin Calculus and Related Topics}, 2nd ed.
Springer, Berlin, 2006.

\bibitem{Rudin-FA}
W.~Rudin,
\emph{Functional Analysis}, 2nd ed.
McGraw-Hill, New York, 1991.

\bibitem{SchrammSmirnov2011}
O.~Schramm and S.~Smirnov,
On the scaling limits of planar percolation.
\emph{Ann. Probab.} \textbf{39} (2011), no.~5, 1768--1814.

\bibitem{SSS2017}
E.~Schertzer, R.~Sun and J.~M.~Swart,
The Brownian web, the Brownian net, and their universality.
In: \emph{Advances in Disordered Systems, Random Processes and Some Applications}, P.~Contucci and C.~Giardin\`a (eds.), Cambridge University Press, 2017.

\bibitem{Tsirelson2004}
B.~Tsirelson,
Nonclassical stochastic flows and continuous products.
\emph{Probab. Surv.} \textbf{1} (2004), 173--298.

\end{thebibliography}

\end{document}
